% Source: https://artofproblemsolving.com/community/c3h2794113_bamo8_problem_set
% Stable identity: aops-topic-2794113-q04
Mod $2$: if $d$ were odd, the six terms would alternate in parity, giving three even terms, at most one of which can be $2$. So $2\mid d$.

Mod $3$: if $3\nmid d$, then $p,p+d,p+2d$ are distinct mod $3$, so one of them is divisible by $3$ and equals $3$. It cannot be $p+d$ or $p+2d$ (these exceed $3$ since $d\ge2$), so $p=3$; but then $p+3d$ is a multiple of $3$ larger than $3$. So $3\mid d$.

Mod $5$: if $5\nmid d$, then $p,\dots,p+4d$ are distinct mod $5$, so one equals $5$. As $d\ge6$, only $p$ can be $5$; then $p+5d$ is a multiple of $5$ larger than $5$. So $5\mid d$. $\blacksquare$
